\chapter{Corporate Stocks}

Corporate stock is at once a legal invention, a way of organizing people, and a claim on future distributions. The lecture begins from the simplest image of a company divided into shares, moves outward to the scale of world stock markets, then narrows again to boards, dividends, dilution, financing, and valuation. What follows keeps that progression. We begin with the corporation as a working social technology, and only then ask what a share is worth, why dividends matter, and why a public company can be worth much more than its book equity.

\section{An Organizational Invention and Its Scale}

The lecture begins with a reminder that the corporation is not merely a legal abstraction. It is a way of letting people work together while holding transferable claims on the venture. Shiller's own founding story makes the point in its simplest form. A businessman contributes money, others contribute time, and the immediate problem is not an abstract theorem but a practical one: how do we divide the company into shares when the contributions are different in kind? The answer in the anecdote is not a finely tuned valuation exercise. It is a social settlement backed by a charter and later adjusted by bonuses, that is, by issuing more shares when resentment develops.

That opening matters because it explains why the corporate form spread so widely. The charter anticipates conflict; the board can alter compensation; shares let people enter, leave, and remain partly committed without freezing the organization. The corporation is therefore presented first as an institutional device for combining labor, capital, and incentives.

From that personal story the lecture turns outward to the scale of listed equity in the world economy. Here the useful comparison is not only the dollar value of the stock market, but also the value of listed companies relative to GDP. The slide is worth preserving because it shows both the absolute and relative measures in the form in which the lecture discusses them.

\begin{figure}[h]
\centering
\includegraphics[width=0.9\textwidth]{lecture_09_figure_02.png}
\caption{World Bank stock-market capitalization table.}
\end{figure}

The key ratio is
\begin{equation}
\frac{\text{market capitalization}}{\text{GDP}}.
\end{equation}
The lecture is careful here: this is a stock divided by a flow. We are taking a stock of wealth---the market value of listed corporations---and comparing it with one year's GDP. The comparison is still informative, but only if we remember that it is dimensionally a stock-to-flow ratio.

On the visible slide, the United States appears at roughly
\[
\left(\frac{\text{market cap}}{\text{GDP}},\ \text{US\$ trillions}\right)=(0.817,\ 11.737),
\]
while the United Kingdom appears at
\[
(0.696,\ 1.851).
\]
The lecturer's point is that the United States has both the largest stock market in the world and one of the highest market-capitalization-to-GDP ratios among major countries. He treats that as a historical fact about capitalist development, connected in the lecture to the early spread of corporate law in the United States and to the country's long financial history.

At the same time, he immediately refuses to overstate the importance of stocks. Rounding U.S. market capitalization to \$12 trillion in 2008 and dividing by about 300 million people gives
\begin{equation}
\frac{12\ \text{trillion dollars}}{300\ \text{million people}}
\approx \$40{,}000 \text{ per person}.
\end{equation}
For a family of four this becomes about \$160,000. That is substantial, but it is not the whole economy, and it is not a plausible once-and-for-all solution to poverty. The lecture strengthens the point by comparing stocks with housing wealth and with total household assets more broadly. Corporate stock matters enormously as an organizational form; it is not the only large asset class, and the economy is not exhausted by it.

\section{Corporation, Equity, and Governance}

After the scale of listed equity has been established, the lecture rebuilds the corporation from the inside. A corporation is a legal person. That is the basic conceptual move. It is an artificial person owned by shareholders, with rights and obligations defined by law. The historical aside about Roman \emph{publicani} makes the point that corporate forms are old, but their mass importance is relatively modern.

The lecture then emphasizes the equal treatment of ordinary shares. This is where the term \emph{equity} enters in its central financial meaning. The etymological discussion is informal, but the institutional point is sharp: each share is ordinarily treated on equal terms with the others. In the usual case, one share carries one vote. There are exceptions, including non-voting shares, but the normal picture is a large set of equal claims whose owners elect a board of directors.

The governance structure then unfolds in stages:
\begin{itemize}
\item shareholders elect the board;
\item the board hires and monitors the chief executive officer;
\item the CEO runs the company as an employee of the corporation;
\item directors owe duties of loyalty to the shareholders.
\end{itemize}
This is the lecture's version of shareholder democracy. It is not one-person-one-vote. It is one-share-one-vote. The distinction matters because it makes control proportional to ownership claims.

The board is not described as a ceremonial body. It is a small group, typically informed about the company in detail, meeting regularly, and able in principle to remove management. The Carl Icahn example is used to dramatize this point: shareholder activism is one way in which the nominal structure of shareholder control becomes real. The example is not there to endorse one activist style in particular. It is there to show that the board is not merely decorative.

The lecture also pauses over the U.S. habit of incorporating in a particular state. Corporate law is not completely national in origin. Delaware appears because the choice of legal domicile affects governance rules and taxes. Again the structure is concrete before it is theoretical.

The for-profit versus nonprofit distinction is introduced at exactly this stage because Yale itself serves as the nearby example. A nonprofit corporation has a board and legal personhood, but no shareholders entitled to profits. The lecture's point is that the organizational sociology may look surprisingly similar across the two forms even though the ultimate claimants are different. That contrast sharpens what is special about the for-profit corporation: profits are ultimately meant to leave the firm for the owners.

\section{Shares, Market Capitalization, and Dividends}

Once the corporate structure is in place, the lecture turns from organization to claims. What exactly is a share? The answer begins with simple arithmetic. If the total equity value of the company is $V$ and there are $N$ shares outstanding, then the value per share is
\begin{equation}
P=\frac{V}{N},
\end{equation}
and the market capitalization is
\begin{equation}
\text{Market cap}=PN.
\end{equation}
A shareholder who owns $n$ shares therefore owns the fraction
\begin{equation}
\frac{n}{N}
\end{equation}
of the common equity claim.

The lecture insists on this arithmetic because raw share counts are meaningless without the denominator. If a company has 1,000,000 shares and we own 1,000 shares, then we own
\[
\frac{1000}{1{,}000{,}000}=\frac{1}{1000}
\]
of the equity. If instead the company has 10,000,000 shares, then the same 1,000 shares amount only to
\[
\frac{1000}{10{,}000{,}000}=\frac{1}{10{,}000}.
\]
The transcript garbles the second line, but the intended arithmetic is clear.

Equal treatment then follows. If the company pays money to one common shareholder, it must pay money to all common shareholders in proportion to their holdings. That is why dividends are paid per share.

\subsection{Question \& Answer}

\paragraph{Question.} If the board of directors is not legally forced to pay dividends, why do common shares have value at all?

\paragraph{Answer.} The lecture's answer is that a for-profit company is ultimately a device for distributing profits to its owners, even if the distribution is delayed. A young firm may retain earnings for many years. Microsoft is the example given: it went for a long period without paying dividends because shareholders expected the retained cash to fund future growth. But a company that would never pay out could not justify its stock price as a for-profit enterprise. In principle, the value of a share is tied to expected future dividends. If a board tries indefinitely to hoard profits against the interests of shareholders, the shareholders can replace the board.

As a theoretical benchmark, the lecture therefore points toward the familiar present-value formula
\begin{equation}
P_t=\sum_{j\ge 1}\frac{\mathbb{E}_t[\mathrm{Div}_{t+j}]}{(1+r)^j}.
\end{equation}
The lecture does not derive this formula formally; it uses it to correct a common misconception. We do not buy stocks merely because we hope someone else will later pay a higher price. In theory, a higher price is justified only by a higher expected stream of future dividends.

This logic immediately explains the ex-dividend price adjustment. When the company pays a dividend, it has less cash, and therefore less value. If the dividend per share is $d_t$, then
\begin{equation}
P_t^{\mathrm{ex}}\approx P_t^{\mathrm{cum}}-d_t.
\end{equation}
The lecture treats this as a routine and largely mechanical adjustment, which is why the old stock tables marked ex-dividend dates explicitly. The decline in price is not bad news about the firm. It is just the arithmetic consequence of cash leaving the company.

This is also why ``selling dividends'' is misleading. Buying just before the dividend and buying just after it should come to nearly the same thing, apart from taxes or other frictions. One gets the dividend and pays a higher price, or buys more cheaply and does not get the dividend. The lecture is using a small piece of price arithmetic to discipline a common retail misunderstanding.

\section{Saleable Shares, Financing, Leverage, and the Stock Market}

After dividends, the lecture returns to what makes the corporation such an effective organizational form: shares are saleable. A founder can leave, a new investor can enter, and the company can issue new shares if it needs money for expansion. This is first presented in the small-company setting because the mechanism is easiest to see there. A factory must be built; a broker finds buyers for newly issued shares; the firm gets money without committing to a fixed repayment schedule.

This is equity financing. Debt financing is the alternative. The company can borrow from a bank or issue corporate bonds. Debt promises fixed payments. Equity promises a residual claim on profits. The lecture is explicit that debt changes the position of shareholders because shareholders are paid only after creditors are paid. In symbols, the common-equity claim is the residual
\begin{equation}
\mathrm{SE}=A-L,
\end{equation}
where $A$ denotes assets and $L$ liabilities. This is the core leverage idea in the lecture. Once debt is introduced, equity becomes a thinner and riskier residual slice.

The lecture now asks whether the small-company story still describes very large public firms. This is the natural point at which the whole theory could collapse. If big firms no longer issue much equity, perhaps the stock market is mostly a place where existing claims are traded back and forth without financing real expansion.

\subsection{Question \& Answer}

\paragraph{Question.} If big companies can finance themselves mainly from retained earnings, what is the stock market still doing for capitalism?

\paragraph{Answer.} The lecture first states the challenge in its strongest form. It introduces Marxian skepticism and then gives the more modern version in Meyers's pecking-order theory. On the aggregate U.S. data he cites, roughly 62\% of corporate financing came from retained earnings, only about 6\% from net equity issuance, and the rest from debt. Read one way, that makes equity issuance look secondary.

But the lecture then refuses to stop with the aggregate number. Fama and French answer that net issuance is not the same as gross issuance. If one firm issues shares while another repurchases them, the aggregate net number can be small even though equity markets are active. That is why the lecture also emphasizes the fraction of firms that issue new shares in a year. The conclusion is not that equity finance is absent, but that one must distinguish between net issuance and the wider role of the stock market in financing and, equally important, in price discovery.

The relevant identity is
\begin{equation}
\text{net equity issuance}
=
\text{new shares sold}
-
\text{share repurchases}.
\end{equation}
Once this distinction is made, the stock market reappears not as a static museum of old shares, but as a live mechanism through which claims are priced, issued, retired, and compared.

\section{Dilution, Repurchases, Common and Preferred Stock}

The lecture next returns from aggregate arguments to shareholder arithmetic. Here the key term is \emph{dilution}. Suppose the company begins with 1,000,000 shares outstanding and we own 1,000 shares. We own one-thousandth of the company. Suppose now that the board issues another 1,000,000 shares at \$10 per share. Then our ownership fraction falls to one two-thousandth.

\paragraph{Worked example.}
\begin{align}
N &= 1{,}000{,}000, & n &= 1000, & \frac{n}{N} &= \frac{1}{1000},\\
N' &= 2{,}000{,}000, & n &= 1000, & \frac{n}{N'} &= \frac{1}{2000}.
\end{align}
So our fraction of the company is cut in half. But the company has also received \$10 million in new cash:
\begin{equation}
1{,}000{,}000 \times \$10 = \$10{,}000{,}000.
\end{equation}
That is why the lecture insists that dilution is not automatically bad. It depends on what price the company gets for the new shares and what it does with the money.

\subsection{Question \& Answer}

\paragraph{Question.} Why is dilution not automatically bad, and why is a buyback not automatically good?

\paragraph{Answer.} Because neither event can be judged from the change in share count alone. An equity issue reduces our ownership fraction, but it may also bring valuable new assets into the firm. A repurchase reduces the share count, but it also takes cash out of the firm. What matters is the price at which the company issues or buys back shares and the use to which the cash is put. The lecture therefore treats both dilution and repurchase as context-dependent.

The lecture then shows the reverse side of dilution. If the company buys back shares, the share count falls. If the buyback is conducted proportionately and we tender proportionately, then our ownership fraction can remain unchanged even as we receive cash. The lecture treats this as close in economic effect to a dividend, though tax treatment may differ.

This leads directly into a comparison between common and preferred stock. Common stock is the true residual ownership claim. Preferred stock is safer and more bond-like. It usually promises a fixed dividend, typically lacks ordinary voting rights, and must generally be paid before common dividends can resume. The lecture is careful, though, not to identify preferred stock with a bond outright. If the company does badly enough, preferred dividends can be missed. In the familiar cumulative form, those missed dividends must be made up before common shareholders can again be paid.

The General Motors example serves a political as well as financial purpose. The U.S. government bought preferred rather than common shares in order to supply capital without taking ordinary voting control and thereby seeming to nationalize the firm. The example locates preferred stock exactly where the lecture wants it: between debt and common equity, safer than the latter, but not the ultimate ownership claim.

\section{Dividend Policy and the Lintner Model}

Only after boards, payouts, financing, dilution, and claim priority have been laid out does the lecture ask what determines actual dividend policy. This is where John Lintner enters. The route is important. We do not begin with a model. We begin with interviews. Lintner asked managers and directors why they set dividends as they do, and the lecture reports two main themes.

First, firms think mature companies ought to pay dividends. Second, firms dislike cutting dividends once established. That second point is crucial because it turns an institutional habit into a mathematical form. If firms had a target payout ratio but were indifferent about cuts, the dividend could jump immediately to its target. Because firms fear the signal sent by a later cut, they adjust only partially.

The blackboard frame from the lecture preserves exactly the moment at which the adjustment equation is being written.

\begin{figure}[h]
\centering
\includegraphics[width=0.78\textwidth]{lecture_09_figure_03.png}
\caption{Lintner dividend adjustment on the board.}
\end{figure}

The screenshot does not fully reveal the right-hand side, but the transcript makes the standard form clear. The lecture's canonical statement is
\begin{equation}
\mathrm{Div}_t-\mathrm{Div}_{t-1}
=
\rho\bigl(\tau\,\mathrm{EPS}_t-\mathrm{Div}_{t-1}\bigr)
+
\varepsilon_t,
\end{equation}
where $0<\tau<1$ is the target payout ratio and $0<\rho<1$ is the speed of adjustment.

The logic is transparent once the pieces are named. The target dividend at time $t$ is roughly $\tau\,\mathrm{EPS}_t$. But actual dividends do not jump instantly to target. Instead, they move partway toward it. If current earnings rise sharply, the board does not want to commit at once to a much higher dividend and then face the embarrassment of cutting it later. So the actual change in dividends is proportional only to the gap between the target payout and last period's dividend.

This is the clearest mathematical core of the lecture. It takes a messy institutional fact---boards want to pay, but hate cutting---and writes it as a stable partial-adjustment rule. The lecture also notes that payout ratios appear to have declined over long periods of time. Earlier shareholders were less patient and demanded more direct cash; modern investors have more willingness to let firms retain earnings for internal growth.

\section{Xerox and Microsoft: Book Equity and Market Value}

The lecture closes by returning from general structure to named companies. Xerox is the first case because it illustrates both technological glory and technological displacement. It began as a great growth company, then ran into trouble when digital copying replaced the older xerographic business. Microsoft is the contrasting case: a company still very valuable as an ongoing concern, with a large stock of financial assets and a market capitalization far above book shareholder equity.

The balance-sheet lesson is straightforward. Assets and liabilities are one side of valuation; market cap is another. Book shareholder equity is the residual after subtracting liabilities from assets:
\[
\mathrm{SE}=A-L.
\]
In the lecture, this is interpreted as a liquidation-style book value: what might remain for shareholders if the firm's assets were sold and debts repaid. But the lecture immediately warns that book values are estimates. Buildings and equipment are not continuously auctioned, and some liabilities, such as taxes or unearned revenue, contain timing and accounting elements that complicate the picture.

For Xerox in 2010, the lecture gives roughly 1.4 billion shares outstanding and a stock price of about \$11.80 per share. That gives a market capitalization of about
\begin{align}
\text{Xerox market cap}
&\approx (1.4\ \text{billion})(\$11.80) \\
&\approx \$16.4\text{--}16.5\ \text{billion}.
\end{align}
Against this it places shareholder equity on the books of roughly \$11.9 billion. For Microsoft, the lecture gives about 8.497 billion shares outstanding at about \$26 per share, hence
\begin{align}
\text{Microsoft market cap}
&\approx (8.497\ \text{billion})(\$26)\\
&\approx \$221\ \text{billion},
\end{align}
against book shareholder equity of roughly \$48 billion.

\subsection{Question \& Answer}

\paragraph{Question.} Why can market capitalization exceed shareholder equity on the books?

\paragraph{Answer.} Because a public company is not purchased merely to be liquidated. If investors are buying a going concern, they are buying future profits, organizational capital, reputational capital, and growth opportunities that do not sit on the balance sheet in simple liquidation form. Xerox, in the lecture's reading, is worth more than its book equity because the market does not believe it is simply a pile of assets to be sold off. Microsoft is worth vastly more than its book equity because its expected future earnings as an operating firm are far larger than the residual value of its recorded assets minus liabilities.

The lecture also states the converse. If a firm's stock price fell far enough below its book shareholder equity, that could create pressure for liquidation or takeover. Yet even here the comparison is not exact, because book shareholder equity is itself an accounting estimate, not a continuously observed market liquidation value.

This is the final valuation lesson of the chapter. Market capitalization and book shareholder equity are related, but they answer different questions. One asks what the market will pay for the common equity today. The other asks what accounting statements imply would remain after subtracting liabilities from assets. The gap between the two is where the lecture places the economic meaning of a going concern.

\section{Summary}

We can now see the lecture's structure as a single argument. Corporate stock begins as an organizational solution to the problem of combining people, money, effort, and disagreement. It becomes a global measure of capitalist development, though never the whole economy. Once the corporation is treated as a legal person with a board and shareholders, a share becomes a residual claim whose value is tied, in principle, to expected future dividends.

From there the lecture moves through the mechanics of public equity: saleability, new issuance, leverage, dilution, repurchases, and the distinction between common and preferred stock. The stock market matters not only because shares trade, but because claims are continually priced, issued, and retired. Lintner's model then condenses dividend policy into a partial-adjustment equation. Finally, Xerox and Microsoft show us why book equity, market capitalization, and going-concern value do not coincide. The corporation, in the end, is neither a mere legal shell nor a pure accounting identity. It is an operating structure whose claims are valued through expectations about future distributions.